\documentclass{leveldoc}
\DocUses{foundations,first_bend,slingshot}
\DocFingerprint{30908b90c2beabdd}

\begin{document}

\section{What the level asks}\label{sec:asks}

The game's text:
\leveldescription

An alpha particle ($q = \val{q}$, $m = \val{m}$) is launched from $(0, 11)$ along $+x$ with
kinetic energy $T_0 = \val{Tzero}$. A fixed nucleus $Q_\mathrm{N} = \val{Qn}$ sits at
$\val{rN}$. The detector is the rectangle $0 \le x \le 1$, $18 \le y \le 20$, with no
condition on direction or energy: the first entry solves the shot. The player places up to
three charges, of either sign and of the magnitudes in the level file, inside the rectangle
$2 \le x \le 10$, $2 \le y \le 18$. The flight ends when the alpha touches a charge, leaves
the world or runs out of time (\xsecref{foundations}{sec:board}); $t_{\max} = \val{tmax}$.

The reference solution places one charge, $Q_\mathrm{S} = \val{Qs}$, at $\val{rS}$. It is
attractive, and it sits below the launch, so it pulls the alpha down onto a close pass of the
nucleus. The nucleus turns the alpha back up toward the left, and the alpha enters the
detector through its right face. The formula in the level's text is the Newtonian scattering
angle of one charge for a particle that comes from infinity; it is derived in
§\ref{sec:formula} and it says how close that pass must be. The flight itself has two
centres, so the arrival is computed numerically, by the method of \levelref{slingshot}.

\section{Setup}\label{sec:setup}

The level file, field by field (the appendix reproduces it):
\setuptable

In symbols: $c = \val{c}$, the alpha has radius $a = 0$, and each charge has radius
$R = \val{R}$. The couplings are $\kappa_\mathrm{N} = qQ_\mathrm{N} = \val{kN}$ and
$\kappa_\mathrm{S} = qQ_\mathrm{S} = \val{kS}$. The launch direction in the file is
$(\val{dirX}, 0)$. The world is $-1 \le x \le 31$, $-1 \le y \le 21$.

\begin{figure}[htbp]
\centering
\begin{tikzpicture}
\begin{axis}[axis equal image, width=0.92\linewidth, xmin=-1.6, xmax=31.6, ymin=-1.6, ymax=21.6,
  xtick={0,5,...,30}, ytick={0,5,...,20}, tick label style={font=\small},
  xlabel={$x$}, ylabel={$y$}, label style={font=\small}]
  \draw[black!8] (0,0) grid (30,20);
  \draw[black!45, dashed] (-1,-1) rectangle (31,21);
  \fill[black!6] (2,2) rectangle (10,18);
  \draw[black!30] (2,2) rectangle (10,18);
  \fill[green!50!black, fill opacity=0.25] (0,18) rectangle (1,20);
  \draw[green!40!black] (0,18) rectangle (1,20);
  \addplot[thin, dashed, black!55] table {sphere.dat};
  \addplot[thin, black!45] table {pathbare.dat};
  \addplot[thick, orange!80!black] table {pathweak.dat};
  \addplot[very thick, blue!70!black] table {path.dat};
  \fill[red!70!black] (18,10) circle[radius=0.28];
  \node[font=\small, red!60!black, anchor=south west] at (18.4,10.15) {$Q_\mathrm{N}$};
  \fill[blue!75!black] (2,2) circle[radius=0.28];
  \node[font=\small, blue!60!black, anchor=south west] at (2.4,2.15) {$Q_\mathrm{S}$};
  \draw[-{Stealth[length=6pt]}, thick] (0,11) -- (2.2,11);
  \fill (0,11) circle[radius=0.12];
  \node[font=\small, anchor=south] at (1.1,11.15) {launch};
  \node[font=\small, green!30!black, anchor=west] at (1.3,19) {detector};
\end{axis}
\end{tikzpicture}
\caption{The reference flight (blue) from the launch to the detector, the same launch with a
weaker steering charge (orange; §\ref{sec:others}) and with the nucleus alone (grey). The
dashed circle, of radius $\val{Rsphere}$ about the nucleus, is where §\ref{sec:circle} reads
the encounter off. The pale rectangle is where the player may place charges.}
\label{fig:board}
\end{figure}

\section{Principles and model}\label{sec:model}

\begin{principles}
  \item The field of a charged sphere and the force on a non-polarizable particle,
    \xeqref{foundations}{eq:coulomb} and \xeqref{foundations}{eq:sphereforce}, superposed for
    two charges.
  \item The relativistic equation of motion \xeqref{foundations}{eq:eom}, the launch
    \xeqref{foundations}{eq:launch}, and the conservation of energy
    \xeqref{foundations}{eq:energy}. With one charge, the orbit
    \xeqref{first_bend}{eq:general}.
  \item The board's rules (\xsecref{foundations}{sec:board}).
  \item The classical Runge--Kutta method, cubic Hermite interpolation and the error estimate
    by halving the step (\xsecref{foundations}{sec:numerics}, \xsecref{slingshot}{sec:method}).
\end{principles}

With the nucleus and the steering charge the state $(\vect{x}, \vect{p})$ obeys
\begin{equation}\label{eq:motion}
  \deriv{\vect{x}}{t} = \frac{\vect{p}}{\gamma m}, \qquad
  \deriv{\vect{p}}{t}
    = \kappa_\mathrm{N}\,\frac{\vect{x}-\vect{r}_\mathrm{N}}{\abs{\vect{x}-\vect{r}_\mathrm{N}}^3}
    + \kappa_\mathrm{S}\,\frac{\vect{x}-\vect{r}_\mathrm{S}}{\abs{\vect{x}-\vect{r}_\mathrm{S}}^3},
  \qquad \gamma = \sqrt{1 + \frac{\abs{\vect{p}}^2}{m^2c^2}},
\end{equation}
from $\vect{x}_0 = (0, 11)$ and $\vect{p}_0 = \abs{\vect{p}_0}\,(1, 0)$, until the first
event. The motion stays in the plane (\xsecref{foundations}{sec:plane}).

\begin{assumption}{The charges are held in place}\label{as:fixed}
Neither charge recoils. That is a rule of the game, and it is the idealization the
Geiger--Marsden experiment itself uses for a nucleus much heavier than an alpha particle.
\end{assumption}

\begin{assumption}{No radiation}\label{as:radiation}
The radiated energy \xeqref{foundations}{eq:larmor} is not taken from the alpha. Along the
reference flight it is $\val{Erad}$, a fraction $\val{EradFrac}$ of $T_0$
(§\ref{sec:hand}). The game makes the same approximation and reports the same integral
(\physref{8}).
\end{assumption}

\begin{assumption}{A point particle}\label{as:point}
$a = 0$, and the alpha does not polarize, so \xeqref{foundations}{eq:sphereforce} is exact.
The closest approach to the nucleus is $\val{rMin}$, far outside the sphere of radius $R$.
\end{assumption}

\begin{assumption}{Where the scattering formula applies}\label{as:formula}
The formula \eqref{eq:chi-tan} is Newtonian, has a single charge, and takes the kinetic
energy and the impact parameter at infinity. On this board $c$ is finite, a second charge
is present, and the alpha is launched a finite distance away. Section~\ref{sec:circle}
measures each of those three: relativity moves the one-centre scattering angle by
$\val{chiRelGap}^\circ$, the second charge moves the turn inside the circle by
$\val{turnGap}^\circ$, and reading $T$ and $b$ at the launch instead of at infinity moves
the unsteered angle by $\val{chiReadGap}^\circ$ (§\ref{sec:unsteered}). The arrival
compared with the game is the integral of \eqref{eq:motion}, which assumes none of this.
\end{assumption}

\section{Derivation}\label{sec:derivation}

\subsection{The scattering angle of one charge}\label{sec:formula}

In the Newtonian limit $c \to \infty$, Level~1's orbit equation
\xeqref{first_bend}{eq:orbit} becomes Binet's equation $u'' + u = -\kappa m/L^2$
(\Goldstein{§3.5}), with $u = 1/r$. Consider a repulsive charge, $\kappa > 0$, and a
particle that comes from infinity with kinetic energy $T$ and impact parameter $b > 0$:
the perpendicular distance from the charge to the incoming asymptote. Then
$\abs{\vect{p}_\infty} = \sqrt{2mT}$ and $L = \abs{\vect{p}_\infty}\,b$, so
\begin{equation}\label{eq:alpha}
  \alpha = \frac{\kappa m}{L^2} = \frac{\kappa}{2Tb^2}.
\end{equation}
Put the apsis, the point of closest approach, at $\theta = 0$. The orbit with
$u'(0) = 0$ is
\begin{equation}\label{eq:hyperbola}
  u(\theta) = A\cos\theta - \alpha,
\end{equation}
the Newtonian case of \xeqref{first_bend}{eq:general}. It reaches infinity, $u = 0$, at
$\theta = \pm\theta_\infty$ with $\cos\theta_\infty = \alpha/A$, so the radius sweeps
$2\theta_\infty$ from one asymptote to the other. The particle's direction changes by the
supplement of that angle, the scattering angle
\begin{equation}\label{eq:chi-geom}
  \chi = \pi - 2\theta_\infty.
\end{equation}

It remains to find $A$. The Newtonian energy $T = L^2(u'^2 + u^2)/(2m) + \kappa u$ holds
along the whole orbit. With \eqref{eq:hyperbola}, $u' = -A\sin\theta$ and
$u'^2 + u^2 = A^2 + \alpha^2 - 2A\alpha\cos\theta$. Also $L^2/(2m) = Tb^2$, and
\eqref{eq:alpha} gives $\kappa = 2\alpha Tb^2$, so the energy reads
\[
  Tb^2\bigl(A^2 + \alpha^2 - 2A\alpha\cos\theta\bigr) + 2\alpha Tb^2\bigl(A\cos\theta - \alpha\bigr) = T.
\]
The coefficients of $\cos\theta$ cancel, $-2A\alpha + 2A\alpha = 0$, and what is left is
$Tb^2(A^2 - \alpha^2) = T$, hence
\begin{equation}\label{eq:amplitude}
  A^2 = \alpha^2 + \frac{1}{b^2}.
\end{equation}
Then $\sin\theta_\infty = 1/(Ab)$. From \eqref{eq:chi-geom},
$\chi/2 = \pi/2 - \theta_\infty$, so $\sin(\chi/2) = \cos\theta_\infty = \alpha/A$ and
$\cos(\chi/2) = \sin\theta_\infty = 1/(Ab)$, and
\begin{equation}\label{eq:chi-tan}
  \tan\frac{\chi}{2} = \alpha b = \frac{\kappa}{2Tb}.
\end{equation}
This is the formula in the level's text, with $\kappa = qQ$ and with $T$ and $b$ taken at
infinity. Both expressions, \eqref{eq:chi-geom} and \eqref{eq:chi-tan}, are evaluated in
§\ref{sec:hand}; they agree to $\val{chiAlg}$.

A larger deflection needs a smaller impact parameter. The detector sits back on the launch
side and above the launch line, so the alpha has to leave the nucleus going up and to the
left: a scattering angle well past a right angle, which \eqref{eq:chi-tan} gives only for a
small $b$.

\subsection{The unsteered alpha}\label{sec:unsteered}

With the nucleus alone the launch velocity is horizontal and the nucleus lies one cell below
the line $y = 11$, so the perpendicular distance is $b = \val{bLaunch}$. The formula with
the launch energy $T_0$ then gives $\chi = \val{chiNaive}^\circ$. That number uses $T$ and
$b$ at the launch, and the launch is not at infinity: $\kappa_\mathrm{N}/r_0 = \val{kN}$
over $\val{rNZero}$ cells. The energy the one-centre problem carries out to infinity is
$T_\infty = T_0 + \kappa_\mathrm{N}/r_0 = \val{TbareInf}$, and the impact parameter that
belongs to it is $b_\infty = \val{bBareInf}$ (the Newtonian angular momentum $L/\gamma$
divided by $\sqrt{2mT_\infty}$). With those, \eqref{eq:chi-tan} gives
$\chi = \val{chiBareInf}^\circ$.

Level~1's orbit through the same launch state, extended to both asymptotes by solving
$u(\theta) = 0$ (\xeqref{first_bend}{eq:general}; the roots by stepping $10^{-3}$ and
bisecting), sweeps a relativistic scattering angle $\chi = \val{chiBare}^\circ$. The two
differ by $\val{chiBareGap}^\circ$: on this board relativity barely moves the angle. The
flight with the nucleus alone is integrated in §\ref{sec:others}. It leaves through the top
of the world at an angle $\val{bareAng}^\circ$, already close to that asymptote, and it
never enters the detector. Steering has to cut the impact parameter down.

\subsection{Two charges}\label{sec:two}

The steering charge makes the force non-central, so the angular momentum about the nucleus
is not constant and \eqref{eq:hyperbola} does not describe the flight. The relativistic
problem of two fixed centres has no separated solution: the term that prevents it is
\xeqref{slingshot}{eq:noseparation}. The flight is the initial-value problem
\eqref{eq:motion}, solved by the method of \xsecref{slingshot}{sec:method}.

The state $y = (x, y, p_x, p_y)$ is advanced by the classical Runge--Kutta method
\xeqref{foundations}{eq:rk4} with a fixed step $h$. On each step the position is
interpolated by the cubic Hermite polynomial \xeqref{slingshot}{eq:hermite}, sampled at
$64$ points. The first sample past an event boundary starts a bisection. Events are taken
in the board's order: a collision with either charge ($R + a$), then the world's edge, then
the detector. The state at the event is one more step of the method, of length
$t_\mathrm{e} - t_n$; the reported position is the interpolant's. The error is of order
$h^4$, and a run at $h/2$ estimates it as a fifteenth of the difference. The energy
\begin{equation}\label{eq:energy}
  W = \gamma mc^2 + \frac{\kappa_\mathrm{N}}{\abs{\vect{x}-\vect{r}_\mathrm{N}}}
    + \frac{\kappa_\mathrm{S}}{\abs{\vect{x}-\vect{r}_\mathrm{S}}}
\end{equation}
is constant, and $\abs{\vect{p}}$ at the entry computed from it is an independent check.

What the energy does say on its own is how close the alpha can come. Split off the
nucleus,
\begin{equation}\label{eq:channel}
  E_\mathrm{N} = T + \frac{\kappa_\mathrm{N}}{r_\mathrm{N}}
    = \bigl(W - mc^2\bigr) - \frac{\kappa_\mathrm{S}}{r_\mathrm{S}}.
\end{equation}
A head-on alpha with energy $E_\mathrm{N}$ stops where $\kappa_\mathrm{N}/r = E_\mathrm{N}$.
At the launch $E_\mathrm{N} = \val{EnucLaunch}$, above $T_0$ because the steering charge is
attractive there. The distance $\kappa_\mathrm{N}/T_0 = \val{headOn}$ is only the answer
for an alpha that stops head-on with nothing but its launch energy.

\subsection{The encounter, read off a circle}\label{sec:circle}

The sharp turn happens near the nucleus, where the steering charge is comparatively far.
To separate the two, take the circle of radius $R_\mathrm{s} = \val{Rsphere}$ about the
nucleus (Figure~\ref{fig:board}) and record the state where the integrated flight crosses
it, inward and outward, with the same interpolant as an event.

From the inward state, continue with the nucleus alone. That is Level~1's orbit
\xeqref{first_bend}{eq:general} from that state: $u(\theta)$ as written there, with this
state's $r$, $p_r$ and $L_z$. Three angles come out of it.
\begin{itemize}
  \item The change in the direction of $\vect{p}$ between the two crossings of the circle,
    the turn the nucleus alone produces inside it. The outgoing crossing is the root of
    $u(\theta) = 1/R_\mathrm{s}$ after closest approach.
  \item The scattering angle of that same orbit from asymptote to asymptote, the root pair
    of $u(\theta) = 0$. This is the relativistic $\chi$ of the one-centre problem defined
    by the inward state.
  \item The Newtonian $\chi$ of \eqref{eq:chi-tan} with the $T_\infty$ and $b_\infty$ of
    that problem: $T_\infty = T + \kappa_\mathrm{N}/R_\mathrm{s}$, and $b_\infty$ from
    $L/\gamma$ and $\sqrt{2mT_\infty}$.
\end{itemize}
The two-centre turn is the change in the direction of $\vect{p}$ between the same two
crossings of the integrated flight. Section~\ref{sec:hand} evaluates all four. The gap
between the two-centre turn and the one-centre turn is the steering charge's effect during
the encounter; the gap between the relativistic and the Newtonian scattering angles is
relativity. On the circle itself $\kappa_\mathrm{N}/R_\mathrm{s} = \val{potIn}$, a fraction
$\val{potFrac}$ of the local kinetic energy, so $T$ and $b$ read off on the circle are not
the $T$ and $b$ of \eqref{eq:chi-tan}.

\section{The reference solution by hand}\label{sec:hand}

\begin{procedure}
  \item\label{st:launch} \emph{The launch} \xeqref{foundations}{eq:launch}:
    $\gamma_0 = \val{gammaZero}$ and $\abs{\vect{p}_0} = \val{pZero}$, so
    $\vect{p}_0 = (\val{pZero}, 0)$.
  \item\label{st:energy} \emph{The energy} \eqref{eq:energy}: the launch distances are
    $r_\mathrm{N} = \val{rNZero}$ and $r_\mathrm{S} = \val{rSZero}$, so
    $V(\vect{x}_0) = \val{VZero}$, $W = \val{W}$ and $W - mc^2 = \val{Wbind}$.
    By \eqref{eq:channel}, $E_\mathrm{N} = \val{EnucLaunch}$ at the launch.
  \item\label{st:unsteered} \emph{No steering charge.} With $b = \val{bLaunch}$ and
    $T = T_0$, \eqref{eq:chi-tan} gives $\chi = \val{chiNaive}^\circ$. With the asymptotic
    values $T_\infty = \val{TbareInf}$ and $b_\infty = \val{bBareInf}$ it gives
    $\chi = \val{chiBareInf}^\circ$, and Level~1's orbit gives $\chi = \val{chiBare}^\circ$
    (§\ref{sec:unsteered}). The alpha is turned upward, not back into the detector.
  \item\label{st:flight} \emph{The reference flight}: the method of §\ref{sec:two} with
    $h = \val{H}$. The states every two time units, and the drift of \eqref{eq:energy}:
    \begin{center}\small\val{states}\end{center}
    The alpha drifts down as $Q_\mathrm{S}$ pulls on it, then between $t = 12$ and
    $t = 16$ the nucleus turns it around. The closest approach is $r = \val{rMin}$ at
    $t = \val{tMin}$, at $(\val{xMin}, \val{yMin})$, where the kinetic energy has fallen to
    $T = \val{Tmin}$. There $E_\mathrm{N} = \val{EnucMin}$, whose head-on stopping distance
    is $\val{rStop}$: the pass is close to head-on and still misses, since
    $\val{rMin} > R$. The steering charge is then $\val{rSteerMin}$ away, and its force is
    a fraction $\val{forceRatio}$ of the nuclear force. The turn is the nucleus's.
  \item\label{st:entry} \emph{The entry.} The step from $t_n = \val{tBefore}$ starts at
    $(\val{xBefore}, \val{yBefore})$, with momentum $(\val{pxBefore}, \val{pyBefore})$.
    It ends at $(\val{xAfter}, \val{yAfter})$, inside the detector. The interpolant crosses
    the right face $x = 1$ at the fraction $u = \val{uE}$ of the step, so
    $t_\mathrm{e} = \val{tE}$ and $y_\mathrm{e} = \val{yE}$, which lies between $18$ and
    $20$. One step of length $t_\mathrm{e} - t_n$ gives
    $\vect{p}_\mathrm{e} = (\val{pxE}, \val{pyE})$, at $\val{angE}^\circ$ to $+x$ and with
    kinetic energy $\val{TE}$. The energy gives
    $\abs{\vect{p}_\mathrm{e}} = \val{pEnergy}$, and the integration matches it to
    $\val{pDiff}$. The alpha has not met a charge or the world's edge, and
    $t_\mathrm{e} \ll t_{\max}$, so the outcome is arrival.
  \item\label{st:error} \emph{The error.} Halving the step moves the entry by
    $\val{errStep}$, of which the time contributes $\val{errT}$. Both runs stop on the face
    $x = 1$, so the abscissa does not move. Along the flight the energy stays within
    $\val{wDev}$ of $W$. The same arrival, coarser:
    \begin{center}\small\val{convergence}\end{center}
    At $h = 0.04$, $\val{roughSteps}$ steps, the entry is still right to $\val{roughErr}$,
    far inside the cell the detector occupies.
  \item\label{st:circle} \emph{The circle} $R_\mathrm{s} = \val{Rsphere}$. The flight
    crosses inward at $t = \val{tIn}$, at $(\val{xIn}, \val{yIn})$, with momentum
    $(\val{pxIn}, \val{pyIn})$: direction $\val{angIn}^\circ$, kinetic energy
    $T = \val{Tin}$, angular momentum $L = \val{Lin}$ and perpendicular distance
    $b = \val{bIn}$. It crosses outward at $t = \val{tOut}$, at
    $(\val{xOut}, \val{yOut})$, in the direction $\val{angOut}^\circ$. The direction has
    changed by $\val{turnTwo}^\circ$.
  \item\label{st:orbit} \emph{The nucleus alone, from that state.} Level~1's orbit leaves
    the circle after sweeping $\theta = \val{thetaExit}$, in the direction
    $\val{angOne}^\circ$: a turn of $\val{turnOne}^\circ$, which is $\val{turnGap}^\circ$
    less than the two-centre turn. Extended to its asymptotes the same orbit scatters by
    $\chi = \val{chiRel}^\circ$. For it, $\Gamma^2 = \val{gammaTwo}$
    (\xeqref{first_bend}{eq:orbit}, with $\kappa_\mathrm{N}/(cL) = \val{eta}$). Its
    asymptotic energy and impact parameter are $T_\infty = \val{Tinf}$ and
    $b_\infty = \val{bInf}$, and \eqref{eq:chi-tan} with those gives
    $\chi = \val{chiInf}^\circ$, within $\val{chiRelGap}^\circ$ of the relativistic orbit.
    The same formula with the local $T$ and $b$ instead gives $\chi = \val{chiTan}^\circ$
    (\eqref{eq:chi-geom} gives $\val{chiGeom}^\circ$; the two differ by $\val{chiAlg}$).
    That is \ref{as:formula}: the formula, fed with the values it asks for, agrees with
    the one-centre orbit; the second charge, not relativity, is what moves the turn.
  \item\label{st:radiation} \emph{The radiation.} The power
    \xeqref{foundations}{eq:larmor} with the acceleration
    \xeqref{first_bend}{eq:acceleration}, summed over the steps by Simpson's rule
    \xeqref{foundations}{eq:simpson}: $E_\mathrm{rad} = \val{Erad}$. Halving the step
    changes it by the relative amount $\val{radErr}$.
\end{procedure}

\begin{result}[The reference flight]
The alpha enters the detector through its right face at $(1, \val{yE})$ at
$t_\mathrm{e} = \val{tE}$, with momentum $(\val{pxE}, \val{pyE})$. The level is solved.
The nucleus did the turning: inside the circle of radius $\val{Rsphere}$ the direction
changes by $\val{turnTwo}^\circ$, of which the nucleus alone accounts for
$\val{turnOne}^\circ$, against $\val{chiInf}^\circ$ from \eqref{eq:chi-tan}.
\end{result}

By hand, Step~\ref{st:flight} is $\val{steps}$ steps of four evaluations of
\eqref{eq:motion}. The convergence table is the shorter route: $\val{roughSteps}$ steps
already place the entry to $\val{roughErr}$.

\section{Comparison with the game}\label{sec:comparison}

\comparisontable

The step-halving error of the reference entry is $\val{errStep}$, and of the two other
flights $\val{weakErr}$ and $\val{bareErr}$, all far below $10^{-10}$. The differences in
the table are therefore the game's. The largest is $\val{cmpWeakXDiff}$, on the abscissa
where the weaker charge leaves the world; it sits at the accuracy the verified integration
has (\xsecref{foundations}{sec:verification}). The tolerance $10^{-9}$ was set from that
accuracy, and from this method's own error, before the comparison. The radiated energy
agrees to the relative difference $\val{cmpRadDiff}$, against a quadrature uncertainty of
$\val{radErr}$.

\section{Other solutions}\label{sec:others}

Two placements besides the reference, each integrated by the same method and compared above.

\begin{enumerate}
  \item The nucleus alone. The alpha leaves through the top of the world,
    $y = 21$, at $(\val{bareX}, \val{bareY})$ at $t = \val{bareT}$, travelling at
    $\val{bareAng}^\circ$. The closest approach is $\val{bareRmin}$. This is the
    backscatter of §\ref{sec:unsteered}: real, and aimed at the ceiling rather than at the
    detector. A single impact parameter of one cell is not small enough.
  \item $Q = \val{detOneQ}$ at $\val{detOneNode}$, so $\kappa = \val{detOneK}$ instead of
    $\kappa_\mathrm{S} = \val{kS}$. The alpha still reaches $r = \val{weakRmin}$ and still
    comes back, leaving at $\val{weakAng}^\circ$ through the same edge, at
    $(\val{weakX}, \val{weakY})$ at $t = \val{weakT}$. That ray crosses $y = 21$ while
    $x = \val{weakX}$ is still to the right of the detector. The reference charge pulls
    the alpha lower before the encounter, and the outgoing ray passes through the cell
    $18 \le y \le 20$ at $x = 1$.
\end{enumerate}

The reference works because $\kappa_\mathrm{S}$ is attractive and placed below the launch:
it reduces the impact parameter from $\val{bLaunch}$ to the $b_\infty = \val{bInf}$ of the
encounter, and \eqref{eq:chi-tan} then gives $\chi = \val{chiInf}^\circ$. During
the encounter itself the steering force is the small fraction $\val{forceRatio}$ of the
nuclear force, and the turn moves by $\val{turnGap}^\circ$. Any charge that arranges a
pass of about this $b$ and this energy will do the same; the two failures above are passes
that come back too high.

\section{What the game computes}\label{sec:engine}

The game integrates \eqref{eq:motion} with its eighth-order Runge--Kutta method and
step-size control at the tolerances $10^{-10}$ and $10^{-12}$ (\physref{5}). The verified
run of the reference takes $\val{gameSteps}$ steps, where this document takes
$\val{steps}$ steps of a fourth-order method at fixed $h$. It finds each crossing on its
interpolant (\physref{6}) and verifies the outcome by comparing its two runs (\physref{7}).
The model is this document's, including \ref{as:radiation} (\physref{8}). The scattering
formula is not part of what the game computes: the agreement in §\ref{sec:comparison} is
between two integrations of \eqref{eq:motion}.

\section{Further reading}\label{sec:reading}

\begin{itemize}
  \item \levelref{first_bend}: the orbit about one charge, of which \eqref{eq:hyperbola}
    is the Newtonian limit and which §\ref{sec:circle} continues through the encounter.
  \item \levelref{slingshot}: two charges, the absence of a separated solution, and the
    numerical method used here.
  \item \Goldstein{§3.5}: Binet's equation. The scattering angle \eqref{eq:chi-tan} is
    derived from it in §\ref{sec:formula}.
\end{itemize}

\levelfileappendix

\end{document}
